% !TEX root = ../AnonymousSubmission2027.tex
\section{Abstract}

% 接触丰富的灵巧操作要求机器人持续感知多指接触状态，并根据快速变化的物理交互调整动作。然而，现有触觉策略通常将密集触觉信号直接注入策略，或仅利用触觉反馈修正动作；预测式和预测-反应式方法虽然引入了未来触觉，却通常只在动作执行前进行一次性预测。当滑移、接触错位或执行误差发生时，固定的触觉预见可能偏离真实交互，造成预测与控制不一致。为此，我们提出 ReTouch，一种具有递归触觉预见能力的接触引导 VLA。ReTouch 使用 \Encoder，按照手指身份和指尖接触拓扑将密集触觉信号编码为结构化表示，并通过双动作专家学习与动作生成相关的 future tactile latents。在执行过程中，ReTouch 利用最新触觉历史异步更新 future tactile latents 与剩余动作，使触觉预见成为随交互递归维护的控制表示。我们进一步搭建 XHand--UR7e 触觉操作平台，并构建包含 800 条真机示范七个接触丰富任务的 XHT-Bench。ReTouch 在标准和分布外测试中的平均成功率均高于最强基线，消融实验进一步验证了各项设计的有效性。

Fusing tactile signals has become an emerging solution in contact-rich dexterous manipulation tasks, enabling robots to accurately perceive multi-finger contact and rapidly changing physical interactions states. Existing tactile fusion policies typically map current tactile signals directly into flattened input tokens for action generation or refinement, while predictive fusion methods also forecast future tactile signals to improve action prediction. However, the generic tactile tokenizer for a dexterous hand may lose structured multi-finger information, and one-step future tactile estimation can also affect precise hand pose estimation.
To mitigate these issues, we introduce ReTouch, a contact-guided VLA model with recursive tactile foresight. Specifically, ReTouch uses a \Encoder to encode dense tactile signals into structured representations based on finger identity and fingertip contact topology, and dual action experts to learn action-relevant future tactile latents. During execution, ReTouch uses the latest tactile history to asynchronously update these latents and the remaining actions, making tactile foresight a control representation that is recursively maintained throughout interaction. We further build an XHand--UR7e tactile manipulation platform and construct XHT-Bench, comprising 800 real-world demonstrations across seven contact-rich tasks. ReTouch achieves a higher average success rate than the strongest baseline under both standard and out-of-distribution settings, while ablation studies validate the effectiveness of its key components.